\documentclass[11pt]{article}

\usepackage{fullpage}
\usepackage{amsmath, amssymb, bm, cite, epsfig, psfrag}
\usepackage{graphicx}
\usepackage{float}
\usepackage{amsthm}
\usepackage{amsfonts}
\usepackage{cite}
\usepackage{hyperref}
\usepackage{pgf,tikz}
\usepackage{enumitem}
\usepackage{mathtools}
\usepackage{siunitx}
\usepackage{../../styles/course}


\begin{document}

\title{Introduction to Hardware Design\\
Shared Memory:  Figures and Code Snippets}
\author{Profs. Sundeep Rangan, Siddharth Garg}
\date{}

\maketitle
\subsection{AXI-MM queue}

\begin{pythoncode}
class PolyCmdHdr(DataList):
    elements = {
        "read_ind": {
            "schema": IntField.specialize(bitwidth=16, signed=False),
            "description": "Read pointer",
        },
        "write_ind": {
            "schema": IntField.specialize(bitwidth=16, signed=False),
            "description": "Write pointer",
        },
        "queue_len": {
            "schema": IntField.specialize(bitwidth=16, signed=False),
            "description": "Maximum queue length in number of entries"
        }
    }
\end{pythoncode}

\begin{ccode}
void kernel(uint64_t queue_base,
            volatile uint32_t *doorbell,
            ap_uint<mem_dwidth> *mem) {
#pragma HLS INTERFACE s_axilite port=queue_base
#pragma HLS INTERFACE s_axilite port=doorbell
#pragma HLS INTERFACE m_axi     port=mem
#pragma HLS INTERFACE ap_ctrl_none port=return
    while (1) {
        // Poll doorbell
        if (*doorbell != 0) {

            // Clear doorbell
            *doorbell = 0;

            // Read queue entries from memory
            // Process all outstanding commands
            // Write completions back to memory
        }
    }
}

\end{ccode}
\end{document}

  

